\section{Related Work}

\textbf{Fixed graphs and architectural structure.} Deep DLGNs separate gate-function learning from the fixed predecessor graph, enabling Boolean inference after hardening \cite{petersen2022}. Convolutional DLGNs introduce shared logic trees, pooling, and residual initialization to exploit image structure \cite{petersen2024}. LogicConnect assigns predecessors within these architectural constraints, preserving the dense or convolutional operators. Degree-preserving controls distinguish the assigned pairs from their degree distribution.

\textbf{Optimizing the predecessor graph.} Connection learning addresses the same graph constraint through a different training problem. Mommen et al. optimize distributions over subsets of candidate predecessors \cite{mommen2025}. LILogic Net develops sparse Top-K routing and basis projection to improve the accuracy and implementation efficiency of compact logic networks \cite{fojcik2026}. These methods can discover connections that a fixed constructor misses, but must process candidate signals during training. The relevant comparison includes accuracy, peak GPU allocation, trainable parameters, and observed training time.

\textbf{Improving gate optimization.} Connectivity is only one source of training difficulty. Mind the Gap uses Gumbel noise and a straight-through estimator to reduce discretization loss \cite{yousefi2025}. Light DLGNs change the gate parameterization, while WARP uses a Walsh–Hadamard representation of Boolean functions \cite{ruttgers2025,gerlach2026}. Multilinear polynomial fitting further examines the redundant optimization directions in a soft mixture of gates \cite{kim2026}. These approaches change how functions are learned on available inputs. LogicConnect changes which inputs are available, leaving the baseline gate-training rule fixed for attribution. Compatibility with another parameterization remains an empirical question, which we examine through matched adapted-WARP controls.

\textbf{Broader design and deployment choices.} BitLogic treats encoding, connectivity, fan-in, and training as interacting design axes for lookup-table-native networks \cite{buehrer2026}. Output-layer studies likewise show that temperature and aggregation choices affect scalability to more classes \cite{brandle2025}. After training, unit tying reduces circuit size through functional approximation and refinement \cite{lee2026}. These studies separate predecessor design from gate representation, output aggregation, and circuit simplification. Their individual gains cannot be added without evaluating the combined system. We therefore separate matched comparisons from published references and adaptations.
